\documentclass{article}
\usepackage{graphicx}
\usepackage{amsmath}
\usepackage{cite}
\usepackage{color}
\usepackage{enumitem}
\usepackage{hyperref}
\usepackage{natbib}
\usepackage{tabularx}
\usepackage{natbib}
\usepackage{ragged2e}

\title{\textbf{La digitalizzazione nelle PA}}
\author{\texttt{Alessandro Meloni GEPID}}
\date{\textit{Anno accademico 2023/2024}}

\begin{document}
\maketitle
    \includegraphics[width = 0.9\linewidth]{Img.jpg}
\centering \tableofcontents
\newpage\centering
\section{Introduzione}
\flushleft \begin{justify}
Ho deciso di presentare questo argomento, che comunque risulta essere molto ampio, soprattutto con lo sviluppo delle tecnologie odierne, e, la quale può aprire un forte dibattito. Specificatamente l'Anac è un ente alla quale mi sono avvicinato durante il mio periodo di studio in triennale, avendo fatto il tirocinio presso l'Ufficio Anticorruzione del Comune di Cagliari. È un ente molto complesso e che richiede un grosso impegno nel poterlo comprendere, e che sicuramente è uno dei settori che potrebbe essere maggiormente colpito (sia in positivo che in negativo, dipende dalla propria opinione) dall'avvento di tecnologie quali Blockchain, Ai, rispetto ai normali database che vengono utilizzati attualmente. Nello specifico ci soffermeremo sull'art 30 del nuovo codice dei contratti pubblici (d.lgs 36/2023), in materia di utilizzo di meccanismi automatizzati nel ciclo di vita dei contratti pubblici. Con questo paper cerco di porre un'analisi oggettiva su quali potrebbero essere le ripercussioni, a livello organizzativo, nel mondo del lavoro e anche nel rispetto dei principi amministrativi dedotti dall'art 97 Cost e varie leggi specifiche in materia.
\end{justify}

\newpage\centering
\section{Nuovo codice degli appalti}
\flushleft Esempio.

\flushleft \subsection{Digitalizzazione dei contratti e principi}
Esempio.
\flushleft \subsection{Art 30 d.lgs 36/2023}
Esempio.
\flushleft \subsection{Rapporto Blockchain/AI e contratti pubblici}
Esempio.

\newpage \centering
\section{Conclusione e opinioni finali}
\flushleft Esempio.

\end{document}
